```latex
\documentclass[12pt,a4paper]{article}

% ============================================================
% PACOTES
% ============================================================

% Codificação e idioma
\usepackage[utf8]{inputenc}
\usepackage[T1]{fontenc}
\usepackage[brazil]{babel}

% Matemática
\usepackage{amsmath}
\usepackage{amssymb}
\usepackage{amsfonts}
\usepackage{mathtools}

% Figuras
\usepackage{graphicx}
\usepackage{float}
\usepackage{subcaption}

% Tabelas
\usepackage{booktabs}
\usepackage{array}
\usepackage{longtable}
\usepackage{multirow}
\usepackage{tabularx}

% Layout
\usepackage[a4paper,
top=2.5cm,
bottom=2.5cm,
left=2.5cm,
right=2.5cm]{geometry}

% Espaçamento
\usepackage{setspace}
\onehalfspacing

% Links e referências
\usepackage[hidelinks]{hyperref}
\usepackage[nameinlink,noabbrev]{cleveref}

% Listas
\usepackage{enumitem}

% Citações
\usepackage[backend=biber,
style=abnt,
sorting=none]{biblatex}

% Unidades
\usepackage{siunitx}

% Código
\usepackage{xcolor}
\usepackage{listings}

% ============================================================
% CONFIGURAÇÕES
% ============================================================

\graphicspath{{figuras/}}

\setlength{\parindent}{1.25cm}
\setlength{\parskip}{0pt}

% Configuração para código
\lstset{
	basicstyle=\ttfamily\small,
	breaklines=true,
	frame=single,
	columns=fullflexible
}

% ============================================================
% COMANDOS PERSONALIZADOS
% ============================================================

\newcommand{\R}{\mathbb{R}}
\newcommand{\N}{\mathbb{N}}
\newcommand{\X}{\mathbf{x}}
\newcommand{\norm}[1]{\left\lVert #1 \right\rVert}
\newcommand{\abs}[1]{\left\lvert #1 \right\rvert}

% ============================================================
% ARQUIVO DE REFERÊNCIAS
% ============================================================

\addbibresource{referencias.bib}

% ============================================================
% INFORMAÇÕES DO DOCUMENTO
% ============================================================

\title{
	\textbf{Análise Matemática das Funções de Base\\
		do Algoritmo FEMa}
}

\author{
	Rennan Furlaneto Collado
}

\date{\today}

% ============================================================
% DOCUMENTO
% ============================================================

\begin{document}
	
	\maketitle
	
	\tableofcontents
	
	\newpage
	
	% ============================================================
	% 1. INTRODUÇÃO
	% ============================================================
	
	\section{Introdução}
	
	O algoritmo FEMa (\textit{Finite Element Machine}) utiliza funções de
	base para determinar a influência das amostras do conjunto de treinamento
	sobre uma nova amostra. Diferentes funções de base produzem diferentes
	comportamentos em relação à distância entre os pontos, podendo alterar a
	forma como as informações locais e globais do conjunto de dados são
	consideradas pelo modelo.
	
	Este documento apresenta uma análise matemática das 21 funções de base
	implementadas no FEMa. O objetivo é caracterizar matematicamente cada
	função e investigar propriedades que possam auxiliar na interpretação dos
	resultados experimentais obtidos pelo algoritmo.
	
	A análise considera propriedades como domínio, valor na origem,
	positividade, monotonicidade, continuidade, comportamento assintótico,
	suporte, taxa de decaimento e influência dos parâmetros.
	
	Além da análise analítica, são utilizados gráficos para visualizar o
	comportamento das funções e facilitar a comparação entre diferentes bases.
	
	% ============================================================
	% 2. FUNDAMENTAÇÃO MATEMÁTICA
	% ============================================================
	
	\section{Fundamentação Matemática}
	
	\subsection{Representação das amostras}
	
	Considere um conjunto de treinamento definido por
	
	\begin{equation}
		\mathcal{D}
		=
		\left\{
		(\mathbf{x}_i,y_i)
		\right\}_{i=1}^{n},
	\end{equation}
	
	em que
	
	\begin{equation}
		\mathbf{x}_i \in \R^d
	\end{equation}
	
	representa o vetor de características da $i$-ésima amostra e
	$y_i$ representa seu respectivo rótulo.
	
	Para uma nova amostra $\mathbf{x}$, a distância entre $\mathbf{x}$ e
	uma amostra de treinamento $\mathbf{x}_i$ pode ser representada por
	
	\begin{equation}
		r_i =
		\norm{\mathbf{x}-\mathbf{x}_i}.
	\end{equation}
	
	No caso da distância Euclidiana,
	
	\begin{equation}
		r_i =
		\sqrt{
			\sum_{j=1}^{d}
			(x_j-x_{ij})^2
		}.
	\end{equation}
	
	\subsection{Funções de base radiais}
	
	Uma função de base radial pode ser representada genericamente por
	
	\begin{equation}
		\phi(r),
	\end{equation}
	
	em que $r$ representa a distância entre dois pontos.
	
	Consequentemente, o peso associado à $i$-ésima amostra pode ser escrito
	como
	
	\begin{equation}
		w_i(\mathbf{x})
		=
		\phi(r_i).
	\end{equation}
	
	A forma específica de $\phi$ determina como a influência de uma amostra
	varia conforme sua distância em relação ao ponto avaliado.
	
	\subsection{Normalização dos pesos}
	
	Quando aplicável, os pesos podem ser normalizados segundo
	
	\begin{equation}
		\widetilde{w}_i(\mathbf{x})
		=
		\frac{
			w_i(\mathbf{x})
		}{
			\displaystyle\sum_{j=1}^{n}w_j(\mathbf{x})
		}.
	\end{equation}
	
	Nesse caso,
	
	\begin{equation}
		\sum_{i=1}^{n}
		\widetilde{w}_i(\mathbf{x})
		=
		1.
	\end{equation}
	
	A normalização permite interpretar os pesos como contribuições relativas
	das diferentes amostras.
	
	% ============================================================
	% 3. CRITÉRIOS DE ANÁLISE
	% ============================================================
	
	\section{Critérios de Análise}
	
	As funções de base serão analisadas segundo um conjunto de propriedades
	matemáticas.
	
	\subsection{Domínio}
	
	O primeiro aspecto analisado será o domínio da função,
	
	\begin{equation}
		\mathcal{D}_{\phi}
		=
		\left\{
		r \in \R :
		\phi(r) \text{ está definida}
		\right\}.
	\end{equation}
	
	Também serão identificadas possíveis restrições sobre os parâmetros.
	
	\subsection{Valor na origem}
	
	Será analisado o valor
	
	\begin{equation}
		\phi(0).
	\end{equation}
	
	Essa propriedade é importante porque $r=0$ representa o caso em que o
	ponto avaliado coincide com uma amostra de treinamento.
	
	\subsection{Positividade}
	
	Será investigado se
	
	\begin{equation}
		\phi(r)\geq0
	\end{equation}
	
	para todo $r$ pertencente ao domínio.
	
	Funções positivas podem apresentar uma interpretação direta como pesos
	de influência.
	
	\subsection{Monotonicidade}
	
	A monotonicidade será investigada por meio da primeira derivada,
	
	\begin{equation}
		\phi'(r).
	\end{equation}
	
	Quando
	
	\begin{equation}
		\phi'(r)\leq0,
	\end{equation}
	
	a função é não crescente em relação à distância.
	
	\subsection{Continuidade e regularidade}
	
	Será analisada a continuidade de $\phi(r)$ e, quando pertinente, de suas
	derivadas.
	
	Quando possível, será determinada a maior classe de regularidade
	
	\begin{equation}
		\phi \in C^k.
	\end{equation}
	
	\subsection{Comportamento assintótico}
	
	Será analisado o limite
	
	\begin{equation}
		\lim_{r\rightarrow\infty}\phi(r).
	\end{equation}
	
	Esse comportamento permite avaliar a influência de pontos cada vez mais
	distantes.
	
	\subsection{Comportamento local}
	
	O comportamento da função próximo da origem também será investigado.
	
	Quando aplicável, será utilizada a expansão de Taylor
	
	\begin{equation}
		\phi(r)
		=
		\phi(0)
		+
		\phi'(0)r
		+
		\frac{\phi''(0)}{2}r^2
		+
		O(r^3).
	\end{equation}
	
	\subsection{Suporte}
	
	Quando aplicável, será determinado o suporte da função,
	
	\begin{equation}
		\operatorname{supp}(\phi)
		=
		\left\{
		r:
		\phi(r)\neq0
		\right\}.
	\end{equation}
	
	Será verificado especialmente se a função possui suporte compacto.
	
	\subsection{Interpolação}
	
	Será investigado o comportamento da função no ponto $r=0$ e sua
	capacidade de produzir pesos compatíveis com a interpolação dos pontos
	de treinamento.
	
	\subsection{Partição da unidade}
	
	Quando houver normalização dos pesos, será verificada a propriedade
	
	\begin{equation}
		\sum_{i=1}^{n}
		\widetilde{w}_i(\mathbf{x})
		=
		1.
	\end{equation}
	
	% ============================================================
	% 4. VISUALIZAÇÃO
	% ============================================================
	
	\section{Metodologia de Visualização}
	
	Para complementar a análise analítica, serão produzidas representações
	gráficas das funções de base.
	
	A primeira visualização consiste na representação unidimensional
	
	\begin{equation}
		y=\phi(r),
	\end{equation}
	
	em função da distância $r$.
	
	Quando pertinente, também serão utilizadas visualizações bidimensionais
	ou mapas de superfície para representar
	
	\begin{equation}
		z=\phi\left(
		\sqrt{x^2+y^2}
		\right).
	\end{equation}
	
	Os gráficos deverão utilizar os mesmos intervalos de distância sempre
	que a comparação direta entre diferentes funções for desejada.
	
	% ============================================================
	% 5. BASES DE ELEMENTOS FINITOS
	% ============================================================
	
	\section{Análise das Funções de Base}
	
	Esta seção apresenta a análise individual das 21 funções de base.
	
	% ============================================================
	% 5.1 SHEPARD
	% ============================================================
	
	\subsection{Shepard}
	
	\subsubsection{Definição}
	
	A função de base Shepard utilizada no FEMa é definida por
	
	\begin{equation}
		\phi_{\mathrm{Shepard}}(r)
		=
		\boxed{\text{INSERIR EXPRESSÃO}}
	\end{equation}
	
	em que $r$ representa a distância entre o ponto avaliado e a amostra
	de treinamento.
	
	\subsubsection{Parâmetros}
	
	Os parâmetros da função são
	
	\begin{equation}
		\boxed{
			\theta = \text{INSERIR PARÂMETROS}
		}
	\end{equation}
	
	Descrever o significado de cada parâmetro.
	
	\subsubsection{Domínio}
	
	\begin{equation}
		\boxed{
			r\in\text{INSERIR DOMÍNIO}
		}
	\end{equation}
	
	\subsubsection{Valor na origem}
	
	\begin{equation}
		\phi_{\mathrm{Shepard}}(0)
		=
		\boxed{\text{INSERIR RESULTADO}}
	\end{equation}
	
	\subsubsection{Positividade}
	
	\begin{equation}
		\phi_{\mathrm{Shepard}}(r)
		\geq0
	\end{equation}
	
	\textbf{Conclusão:} \textit{INSERIR ANÁLISE.}
	
	\subsubsection{Monotonicidade}
	
	A primeira derivada é
	
	\begin{equation}
		\phi_{\mathrm{Shepard}}'(r)
		=
		\boxed{\text{INSERIR DERIVADA}}
	\end{equation}
	
	\textbf{Conclusão:} \textit{INSERIR ANÁLISE.}
	
	\subsubsection{Comportamento assintótico}
	
	\begin{equation}
		\lim_{r\rightarrow\infty}
		\phi_{\mathrm{Shepard}}(r)
		=
		\boxed{\text{INSERIR RESULTADO}}
	\end{equation}
	
	\subsubsection{Continuidade e regularidade}
	
	\textit{INSERIR ANÁLISE.}
	
	\subsubsection{Suporte}
	
	\textit{INSERIR ANÁLISE.}
	
	\subsubsection{Gráfico}
	
	\begin{figure}[H]
		\centering
		\includegraphics[width=0.85\textwidth]
		{shepard.png}
		\caption{
			Comportamento da função de base Shepard em função da distância.
		}
		\label{fig:shepard}
	\end{figure}
	
	\subsubsection{Observações}
	
	\begin{itemize}
		\item \textit{INSERIR OBSERVAÇÃO SOBRE A REGIÃO PRÓXIMA À ORIGEM.}
		\item \textit{INSERIR OBSERVAÇÃO SOBRE O DECAIMENTO.}
		\item \textit{INSERIR OBSERVAÇÃO SOBRE OS PARÂMETROS.}
		\item \textit{INSERIR OBSERVAÇÃO SOBRE O FEMa.}
	\end{itemize}
	
	% ============================================================
	% 5.2 RADIAL
	% ============================================================
	
	\subsection{Radial}
	
	\subsubsection{Definição}
	
	\begin{equation}
		\phi_{\mathrm{Radial}}(r)
		=
		\boxed{\text{INSERIR EXPRESSÃO}}
	\end{equation}
	
	\subsubsection{Análise matemática}
	
	\begin{itemize}
		\item Domínio: \textit{INSERIR}.
		\item $\phi(0)$: \textit{INSERIR}.
		\item Positividade: \textit{INSERIR}.
		\item Monotonicidade: \textit{INSERIR}.
		\item Continuidade: \textit{INSERIR}.
		\item Comportamento assintótico: \textit{INSERIR}.
		\item Suporte: \textit{INSERIR}.
	\end{itemize}
	
	\subsubsection{Gráfico}
	
	\begin{figure}[H]
		\centering
		\includegraphics[width=0.85\textwidth]
		{radial.png}
		\caption{
			Comportamento da função de base Radial.
		}
		\label{fig:radial}
	\end{figure}
	
	\subsubsection{Observações}
	
	\textit{INSERIR OBSERVAÇÕES.}
	
	% ============================================================
	% 5.3 RBF GAUSSIAN
	% ============================================================
	
	\subsection{RBF Gaussian}
	
	\subsubsection{Definição}
	
	\begin{equation}
		\phi_{\mathrm{Gaussian}}(r)
		=
		\boxed{\text{INSERIR EXPRESSÃO}}
	\end{equation}
	
	\subsubsection{Análise matemática}
	
	\begin{itemize}
		\item Domínio: \textit{INSERIR}.
		\item $\phi(0)$: \textit{INSERIR}.
		\item Positividade: \textit{INSERIR}.
		\item Monotonicidade: \textit{INSERIR}.
		\item Continuidade: \textit{INSERIR}.
		\item Comportamento assintótico: \textit{INSERIR}.
		\item Suporte: \textit{INSERIR}.
	\end{itemize}
	
	\subsubsection{Gráfico}
	
	\begin{figure}[H]
		\centering
		\includegraphics[width=0.85\textwidth]
		{rbf_gaussian.png}
		\caption{
			Comportamento da função RBF Gaussian.
		}
		\label{fig:rbf-gaussian}
	\end{figure}
	
	\subsubsection{Observações}
	
	\textit{INSERIR OBSERVAÇÕES.}
	
	% ============================================================
	% 5.4 MULTIQUADRATIC
	% ============================================================
	
	\subsection{Multiquadratic}
	
	\subsubsection{Definição}
	
	\begin{equation}
		\phi_{\mathrm{Multiquadratic}}(r)
		=
		\boxed{\text{INSERIR EXPRESSÃO}}
	\end{equation}
	
	\subsubsection{Análise}
	
	\begin{itemize}
		\item Domínio: \textit{INSERIR}.
		\item Valor na origem: \textit{INSERIR}.
		\item Positividade: \textit{INSERIR}.
		\item Monotonicidade: \textit{INSERIR}.
		\item Regularidade: \textit{INSERIR}.
		\item Comportamento assintótico: \textit{INSERIR}.
	\end{itemize}
	
	\subsubsection{Gráfico}
	
	\begin{figure}[H]
		\centering
		\includegraphics[width=0.85\textwidth]
		{multiquadratic.png}
		\caption{
			Comportamento da função Multiquadratic.
		}
		\label{fig:multiquadratic}
	\end{figure}
	
	\subsubsection{Observações}
	
	\textit{INSERIR OBSERVAÇÕES.}
	
	% ============================================================
	% 5.5 INVERSE MULTIQUADRATIC
	% ============================================================
	
	\subsection{Inverse Multiquadratic}
	
	\subsubsection{Definição}
	
	\begin{equation}
		\phi_{\mathrm{InverseMultiquadratic}}(r)
		=
		\boxed{\text{INSERIR EXPRESSÃO}}
	\end{equation}
	
	\subsubsection{Análise}
	
	\begin{itemize}
		\item Domínio: \textit{INSERIR}.
		\item Valor na origem: \textit{INSERIR}.
		\item Positividade: \textit{INSERIR}.
		\item Monotonicidade: \textit{INSERIR}.
		\item Regularidade: \textit{INSERIR}.
		\item Comportamento assintótico: \textit{INSERIR}.
	\end{itemize}
	
	\subsubsection{Gráfico}
	
	\begin{figure}[H]
		\centering
		\includegraphics[width=0.85\textwidth]
		{inverse_multiquadratic.png}
		\caption{
			Comportamento da função Inverse Multiquadratic.
		}
		\label{fig:inverse-multiquadratic}
	\end{figure}
	
	\subsubsection{Observações}
	
	\textit{INSERIR OBSERVAÇÕES.}
	
	% ============================================================
	% 5.6 WENDLAND C2
	% ============================================================
	
	\subsection{Wendland C2}
	
	\subsubsection{Definição}
	
	\begin{equation}
		\phi_{\mathrm{Wendland}}(r)
		=
		\boxed{\text{INSERIR EXPRESSÃO}}
	\end{equation}
	
	\subsubsection{Análise}
	
	\begin{itemize}
		\item Domínio: \textit{INSERIR}.
		\item Valor na origem: \textit{INSERIR}.
		\item Positividade: \textit{INSERIR}.
		\item Monotonicidade: \textit{INSERIR}.
		\item Regularidade: \textit{INSERIR}.
		\item Suporte compacto: \textit{SIM/NÃO}.
	\end{itemize}
	
	\subsubsection{Gráfico}
	
	\begin{figure}[H]
		\centering
		\includegraphics[width=0.85\textwidth]
		{wendland_c2.png}
		\caption{
			Comportamento da função Wendland C2.
		}
		\label{fig:wendland-c2}
	\end{figure}
	
	\subsubsection{Observações}
	
	\textit{INSERIR OBSERVAÇÕES.}
	
	% ============================================================
	% 5.7 CUBIC SPLINE
	% ============================================================
	
	\subsection{Cubic Spline}
	
	\subsubsection{Definição}
	
	\begin{equation}
		\phi_{\mathrm{CubicSpline}}(r)
		=
		\boxed{\text{INSERIR EXPRESSÃO}}
	\end{equation}
	
	\subsubsection{Análise}
	
	\begin{itemize}
		\item Domínio: \textit{INSERIR}.
		\item Valor na origem: \textit{INSERIR}.
		\item Positividade: \textit{INSERIR}.
		\item Monotonicidade: \textit{INSERIR}.
		\item Regularidade: \textit{INSERIR}.
		\item Suporte: \textit{INSERIR}.
	\end{itemize}
	
	\subsubsection{Gráfico}
	
	\begin{figure}[H]
		\centering
		\includegraphics[width=0.85\textwidth]
		{cubic_spline.png}
		\caption{
			Comportamento da função Cubic Spline.
		}
		\label{fig:cubic-spline}
	\end{figure}
	
	\subsubsection{Observações}
	
	\textit{INSERIR OBSERVAÇÕES.}
	
	% ============================================================
	% 5.8 QUARTIC SPLINE
	% ============================================================
	
	\subsection{Quartic Spline}
	
	\subsubsection{Definição}
	
	\begin{equation}
		\phi_{\mathrm{QuarticSpline}}(r)
		=
		\boxed{\text{INSERIR EXPRESSÃO}}
	\end{equation}
	
	\subsubsection{Análise}
	
	\begin{itemize}
		\item Domínio: \textit{INSERIR}.
		\item Valor na origem: \textit{INSERIR}.
		\item Positividade: \textit{INSERIR}.
		\item Monotonicidade: \textit{INSERIR}.
		\item Regularidade: \textit{INSERIR}.
		\item Suporte: \textit{INSERIR}.
	\end{itemize}
	
	\subsubsection{Gráfico}
	
	\begin{figure}[H]
		\centering
		\includegraphics[width=0.85\textwidth]
		{quartic_spline.png}
		\caption{
			Comportamento da função Quartic Spline.
		}
		\label{fig:quartic-spline}
	\end{figure}
	
	\subsubsection{Observações}
	
	\textit{INSERIR OBSERVAÇÕES.}
	
	% ============================================================
	% 5.9 GENERALIZED EXPONENTIAL
	% ============================================================
	
	\subsection{Generalized Exponential}
	
	\subsubsection{Definição}
	
	\begin{equation}
		\phi_{\mathrm{GenExp}}(r)
		=
		\boxed{\text{INSERIR EXPRESSÃO}}
	\end{equation}
	
	\subsubsection{Análise}
	
	\begin{itemize}
		\item Domínio: \textit{INSERIR}.
		\item Valor na origem: \textit{INSERIR}.
		\item Positividade: \textit{INSERIR}.
		\item Monotonicidade: \textit{INSERIR}.
		\item Regularidade: \textit{INSERIR}.
		\item Comportamento assintótico: \textit{INSERIR}.
	\end{itemize}
	
	\subsubsection{Gráfico}
	
	\begin{figure}[H]
		\centering
		\includegraphics[width=0.85\textwidth]
		{gen_exponential.png}
		\caption{
			Comportamento da função Generalized Exponential.
		}
		\label{fig:gen-exponential}
	\end{figure}
	
	\subsubsection{Observações}
	
	\textit{INSERIR OBSERVAÇÕES.}
	
	% ============================================================
	% 5.10 SOFTMAX RADIAL
	% ============================================================
	
	\subsection{Softmax Radial}
	
	\subsubsection{Definição}
	
	\begin{equation}
		\phi_{\mathrm{SoftmaxRadial}}(r)
		=
		\boxed{\text{INSERIR EXPRESSÃO}}
	\end{equation}
	
	\subsubsection{Análise}
	
	\begin{itemize}
		\item Domínio: \textit{INSERIR}.
		\item Valor na origem: \textit{INSERIR}.
		\item Positividade: \textit{INSERIR}.
		\item Monotonicidade: \textit{INSERIR}.
		\item Regularidade: \textit{INSERIR}.
		\item Comportamento assintótico: \textit{INSERIR}.
	\end{itemize}
	
	\subsubsection{Gráfico}
	
	\begin{figure}[H]
		\centering
		\includegraphics[width=0.85\textwidth]
		{softmax_radial.png}
		\caption{
			Comportamento da função Softmax Radial.
		}
		\label{fig:softmax-radial}
	\end{figure}
	
	\subsubsection{Observações}
	
	\textit{INSERIR OBSERVAÇÕES.}
	
	% ============================================================
	% 5.11 ATTENTION
	% ============================================================
	
	\subsection{Attention}
	
	\subsubsection{Definição}
	
	\begin{equation}
		\phi_{\mathrm{Attention}}(r)
		=
		\boxed{\text{INSERIR EXPRESSÃO}}
	\end{equation}
	
	\subsubsection{Análise}
	
	\begin{itemize}
		\item Domínio: \textit{INSERIR}.
		\item Valor na origem: \textit{INSERIR}.
		\item Positividade: \textit{INSERIR}.
		\item Monotonicidade: \textit{INSERIR}.
		\item Regularidade: \textit{INSERIR}.
		\item Comportamento assintótico: \textit{INSERIR}.
	\end{itemize}
	
	\subsubsection{Gráfico}
	
	\begin{figure}[H]
		\centering
		\includegraphics[width=0.85\textwidth]
		{attention.png}
		\caption{
			Comportamento da função Attention.
		}
		\label{fig:attention}
	\end{figure}
	
	\subsubsection{Observações}
	
	\textit{INSERIR OBSERVAÇÕES.}
	
	% ============================================================
	% 5.12 LOGARITHMIC
	% ============================================================
	
	\subsection{Logarithmic}
	
	\subsubsection{Definição}
	
	\begin{equation}
		\phi_{\mathrm{Logarithmic}}(r)
		=
		\boxed{\text{INSERIR EXPRESSÃO}}
	\end{equation}
	
	\subsubsection{Análise}
	
	\begin{itemize}
		\item Domínio: \textit{INSERIR}.
		\item Valor na origem: \textit{INSERIR}.
		\item Positividade: \textit{INSERIR}.
		\item Monotonicidade: \textit{INSERIR}.
		\item Regularidade: \textit{INSERIR}.
		\item Comportamento assintótico: \textit{INSERIR}.
	\end{itemize}
	
	\subsubsection{Gráfico}
	
	\begin{figure}[H]
		\centering
		\includegraphics[width=0.85\textwidth]
		{logarithmic.png}
		\caption{
			Comportamento da função Logarithmic.
		}
		\label{fig:logarithmic}
	\end{figure}
	
	\subsubsection{Observações}
	
	\textit{INSERIR OBSERVAÇÕES.}
	
	% ============================================================
	% 5.13 HARMONIC
	% ============================================================
	
	\subsection{Harmonic}
	
	\subsubsection{Definição}
	
	\begin{equation}
		\phi_{\mathrm{Harmonic}}(r)
		=
		\boxed{\text{INSERIR EXPRESSÃO}}
	\end{equation}
	
	\subsubsection{Análise}
	
	\begin{itemize}
		\item Domínio: \textit{INSERIR}.
		\item Valor na origem: \textit{INSERIR}.
		\item Positividade: \textit{INSERIR}.
		\item Monotonicidade: \textit{INSERIR}.
		\item Regularidade: \textit{INSERIR}.
		\item Comportamento assintótico: \textit{INSERIR}.
	\end{itemize}
	
	\subsubsection{Gráfico}
	
	\begin{figure}[H]
		\centering
		\includegraphics[width=0.85\textwidth]
		{harmonic.png}
		\caption{
			Comportamento da função Harmonic.
		}
		\label{fig:harmonic}
	\end{figure}
	
	\subsubsection{Observações}
	
	\textit{INSERIR OBSERVAÇÕES.}
	
	% ============================================================
	% 5.14 LAPLACIAN
	% ============================================================
	
	\subsection{Laplacian}
	
	\subsubsection{Definição}
	
	\begin{equation}
		\phi_{\mathrm{Laplacian}}(r)
		=
		\boxed{\text{INSERIR EXPRESSÃO}}
	\end{equation}
	
	\subsubsection{Análise}
	
	\begin{itemize}
		\item Domínio: \textit{INSERIR}.
		\item Valor na origem: \textit{INSERIR}.
		\item Positividade: \textit{INSERIR}.
		\item Monotonicidade: \textit{INSERIR}.
		\item Regularidade: \textit{INSERIR}.
		\item Comportamento assintótico: \textit{INSERIR}.
	\end{itemize}
	
	\subsubsection{Gráfico}
	
	\begin{figure}[H]
		\centering
		\includegraphics[width=0.85\textwidth]
		{laplacian.png}
		\caption{
			Comportamento da função Laplacian.
		}
		\label{fig:laplacian}
	\end{figure}
	
	\subsubsection{Observações}
	
	\textit{INSERIR OBSERVAÇÕES.}
	
	% ============================================================
	% 5.15 CAUCHY
	% ============================================================
	
	\subsection{Cauchy}
	
	\subsubsection{Definição}
	
	\begin{equation}
		\phi_{\mathrm{Cauchy}}(r)
		=
		\boxed{\text{INSERIR EXPRESSÃO}}
	\end{equation}
	
	\subsubsection{Análise}
	
	\begin{itemize}
		\item Domínio: \textit{INSERIR}.
		\item Valor na origem: \textit{INSERIR}.
		\item Positividade: \textit{INSERIR}.
		\item Monotonicidade: \textit{INSERIR}.
		\item Regularidade: \textit{INSERIR}.
		\item Comportamento assintótico: \textit{INSERIR}.
	\end{itemize}
	
	\subsubsection{Gráfico}
	
	\begin{figure}[H]
		\centering
		\includegraphics[width=0.85\textwidth]
		{cauchy.png}
		\caption{
			Comportamento da função Cauchy.
		}
		\label{fig:cauchy}
	\end{figure}
	
	\subsubsection{Observações}
	
	\textit{INSERIR OBSERVAÇÕES.}
	
	% ============================================================
	% 5.16 STUDENT-T
	% ============================================================
	
	\subsection{Student-t}
	
	\subsubsection{Definição}
	
	\begin{equation}
		\phi_{\mathrm{StudentT}}(r)
		=
		\boxed{\text{INSERIR EXPRESSÃO}}
	\end{equation}
	
	\subsubsection{Análise}
	
	\begin{itemize}
		\item Domínio: \textit{INSERIR}.
		\item Valor na origem: \textit{INSERIR}.
		\item Positividade: \textit{INSERIR}.
		\item Monotonicidade: \textit{INSERIR}.
		\item Regularidade: \textit{INSERIR}.
		\item Comportamento assintótico: \textit{INSERIR}.
	\end{itemize}
	
	\subsubsection{Gráfico}
	
	\begin{figure}[H]
		\centering
		\includegraphics[width=0.85\textwidth]
		{student_t.png}
		\caption{
			Comportamento da função Student-t.
		}
		\label{fig:student-t}
	\end{figure}
	
	\subsubsection{Observações}
	
	\textit{INSERIR OBSERVAÇÕES.}
	
	% ============================================================
	% 5.17 COSINE
	% ============================================================
	
	\subsection{Cosine}
	
	\subsubsection{Definição}
	
	\begin{equation}
		\phi_{\mathrm{Cosine}}(r)
		=
		\boxed{\text{INSERIR EXPRESSÃO}}
	\end{equation}
	
	\subsubsection{Análise}
	
	\begin{itemize}
		\item Domínio: \textit{INSERIR}.
		\item Valor na origem: \textit{INSERIR}.
		\item Positividade: \textit{INSERIR}.
		\item Monotonicidade: \textit{INSERIR}.
		\item Regularidade: \textit{INSERIR}.
		\item Comportamento assintótico: \textit{INSERIR}.
	\end{itemize}
	
	\subsubsection{Gráfico}
	
	\begin{figure}[H]
		\centering
		\includegraphics[width=0.85\textwidth]
		{cosine.png}
		\caption{
			Comportamento da função Cosine.
		}
		\label{fig:cosine}
	\end{figure}
	
	\subsubsection{Observações}
	
	\textit{INSERIR OBSERVAÇÕES.}
	
	% ============================================================
	% 5.18 SIGMOIDAL
	% ============================================================
	
	\subsection{Sigmoidal}
	
	\subsubsection{Definição}
	
	\begin{equation}
		\phi_{\mathrm{Sigmoidal}}(r)
		=
		\boxed{\text{INSERIR EXPRESSÃO}}
	\end{equation}
	
	\subsubsection{Análise}
	
	\begin{itemize}
		\item Domínio: \textit{INSERIR}.
		\item Valor na origem: \textit{INSERIR}.
		\item Positividade: \textit{INSERIR}.
		\item Monotonicidade: \textit{INSERIR}.
		\item Regularidade: \textit{INSERIR}.
		\item Comportamento assintótico: \textit{INSERIR}.
	\end{itemize}
	
	\subsubsection{Gráfico}
	
	\begin{figure}[H]
		\centering
		\includegraphics[width=0.85\textwidth]
		{sigmoidal.png}
		\caption{
			Comportamento da função Sigmoidal.
		}
		\label{fig:sigmoidal}
	\end{figure}
	
	\subsubsection{Observações}
	
	\textit{INSERIR OBSERVAÇÕES.}
	
	% ============================================================
	% 5.19 LORENTZIAN
	% ============================================================
	
	\subsection{Lorentzian}
	
	\subsubsection{Definição}
	
	\begin{equation}
		\phi_{\mathrm{Lorentzian}}(r)
		=
		\boxed{\text{INSERIR EXPRESSÃO}}
	\end{equation}
	
	\subsubsection{Análise}
	
	\begin{itemize}
		\item Domínio: \textit{INSERIR}.
		\item Valor na origem: \textit{INSERIR}.
		\item Positividade: \textit{INSERIR}.
		\item Monotonicidade: \textit{INSERIR}.
		\item Regularidade: \textit{INSERIR}.
		\item Comportamento assintótico: \textit{INSERIR}.
	\end{itemize}
	
	\subsubsection{Gráfico}
	
	\begin{figure}[H]
		\centering
		\includegraphics[width=0.85\textwidth]
		{lorentzian.png}
		\caption{
			Comportamento da função Lorentzian.
		}
		\label{fig:lorentzian}
	\end{figure}
	
	\subsubsection{Observações}
	
	\textit{INSERIR OBSERVAÇÕES.}
	
	% ============================================================
	% 5.20 ENTROPIC
	% ============================================================
	
	\subsection{Entropic}
	
	\subsubsection{Definição}
	
	\begin{equation}
		\phi_{\mathrm{Entropic}}(r)
		=
		\boxed{\text{INSERIR EXPRESSÃO}}
	\end{equation}
	
	\subsubsection{Análise}
	
	\begin{itemize}
		\item Domínio: \textit{INSERIR}.
		\item Valor na origem: \textit{INSERIR}.
		\item Positividade: \textit{INSERIR}.
		\item Monotonicidade: \textit{INSERIR}.
		\item Regularidade: \textit{INSERIR}.
		\item Comportamento assintótico: \textit{INSERIR}.
	\end{itemize}
	
	\subsubsection{Gráfico}
	
	\begin{figure}[H]
		\centering
		\includegraphics[width=0.85\textwidth]
		{entropic.png}
		\caption{
			Comportamento da função Entropic.
		}
		\label{fig:entropic}
	\end{figure}
	
	\subsubsection{Observações}
	
	\textit{INSERIR OBSERVAÇÕES.}
	
	% ============================================================
	% 5.21 RATIONAL QUADRATIC
	% ============================================================
	
	\subsection{Rational Quadratic}
	
	\subsubsection{Definição}
	
	\begin{equation}
		\phi_{\mathrm{RationalQuadratic}}(r)
		=
		\boxed{\text{INSERIR EXPRESSÃO}}
	\end{equation}
	
	\subsubsection{Análise}
	
	\begin{itemize}
		\item Domínio: \textit{INSERIR}.
		\item Valor na origem: \textit{INSERIR}.
		\item Positividade: \textit{INSERIR}.
		\item Monotonicidade: \textit{INSERIR}.
		\item Regularidade: \textit{INSERIR}.
		\item Comportamento assintótico: \textit{INSERIR}.
	\end{itemize}
	
	\subsubsection{Gráfico}
	
	\begin{figure}[H]
		\centering
		\includegraphics[width=0.85\textwidth]
		{rational_quadratic.png}
		\caption{
			Comportamento da função Rational Quadratic.
		}
		\label{fig:rational-quadratic}
	\end{figure}
	
	\subsubsection{Observações}
	
	\textit{INSERIR OBSERVAÇÕES.}
	
	% ============================================================
	% 6. TABELA COMPARATIVA
	% ============================================================
	
	\section{Comparação das Funções de Base}
	
	A Tabela~\ref{tab:comparacao-bases} apresenta uma síntese das
	propriedades matemáticas identificadas nas 21 funções.
	
	\begin{table}[H]
		\centering
		\small
		\caption{
			Comparação das propriedades matemáticas das funções de base.
		}
		\label{tab:comparacao-bases}
		\begin{tabular}{lcccccc}
			\toprule
			\textbf{Base}
			&
			\textbf{Positiva}
			&
			\textbf{Monotônica}
			&
			$\boldsymbol{\phi(0)}$
			&
			\textbf{Decaimento}
			&
			\textbf{Suporte}
			&
			\textbf{Interpolativa}
			\\
			\midrule
			
			Shepard
			& -- & -- & -- & -- & -- & -- \\
			
			Radial
			& -- & -- & -- & -- & -- & -- \\
			
			RBF Gaussian
			& -- & -- & -- & -- & -- & -- \\
			
			Multiquadratic
			& -- & -- & -- & -- & -- & -- \\
			
			Inverse Multiquadratic
			& -- & -- & -- & -- & -- & -- \\
			
			Wendland C2
			& -- & -- & -- & -- & -- & -- \\
			
			Cubic Spline
			& -- & -- & -- & -- & -- & -- \\
			
			Quartic Spline
			& -- & -- & -- & -- & -- & -- \\
			
			Generalized Exponential
			& -- & -- & -- & -- & -- & -- \\
			
			Softmax Radial
			& -- & -- & -- & -- & -- & -- \\
			
			Attention
			& -- & -- & -- & -- & -- & -- \\
			
			Logarithmic
			& -- & -- & -- & -- & -- & -- \\
			
			Harmonic
			& -- & -- & -- & -- & -- & -- \\
			
			Laplacian
			& -- & -- & -- & -- & -- & -- \\
			
			Cauchy
			& -- & -- & -- & -- & -- & -- \\
			
			Student-t
			& -- & -- & -- & -- & -- & -- \\
			
			Cosine
			& -- & -- & -- & -- & -- & -- \\
			
			Sigmoidal
			& -- & -- & -- & -- & -- & -- \\
			
			Lorentzian
			& -- & -- & -- & -- & -- & -- \\
			
			Entropic
			& -- & -- & -- & -- & -- & -- \\
			
			Rational Quadratic
			& -- & -- & -- & -- & -- & -- \\
			
			\bottomrule
		\end{tabular}
	\end{table}
	
	% ============================================================
	% 7. COMPARAÇÃO GRÁFICA
	% ============================================================
	
	\section{Comparação Gráfica}
	
	\subsection{Comparação das funções}
	
	\begin{figure}[H]
		\centering
		\includegraphics[width=\textwidth]
		{comparacao_21_bases.png}
		\caption{
			Comparação do comportamento das 21 funções de base em função
			da distância.
		}
		\label{fig:comparacao-21}
	\end{figure}
	
	\subsection{Funções de decaimento rápido}
	
	\begin{figure}[H]
		\centering
		\includegraphics[width=0.9\textwidth]
		{decaimento_rapido.png}
		\caption{
			Funções de base caracterizadas por maior concentração da
			influência em distâncias próximas.
		}
		\label{fig:decaimento-rapido}
	\end{figure}
	
	\subsection{Funções de decaimento lento}
	
	\begin{figure}[H]
		\centering
		\includegraphics[width=0.9\textwidth]
		{decaimento_lento.png}
		\caption{
			Funções de base com decaimento mais gradual da influência.
		}
		\label{fig:decaimento-lento}
	\end{figure}
	
	\subsection{Funções com suporte compacto}
	
	\begin{figure}[H]
		\centering
		\includegraphics[width=0.9\textwidth]
		{suporte_compacto.png}
		\caption{
			Funções de base com suporte compacto.
		}
		\label{fig:suporte-compacto}
	\end{figure}
	
	% ============================================================
	% 8. INFLUÊNCIA DOS PARÂMETROS
	% ============================================================
	
	\section{Influência dos Parâmetros}
	
	Para as funções que possuem parâmetros de escala, será analisado como
	a variação desses parâmetros modifica o comportamento da função.
	
	Considere uma função genérica
	
	\begin{equation}
		\phi(r;\theta),
	\end{equation}
	
	em que $\theta$ representa um ou mais parâmetros.
	
	\subsection{Variação da largura}
	
	\begin{figure}[H]
		\centering
		\includegraphics[width=0.9\textwidth]
		{influencia_largura.png}
		\caption{
			Influência do parâmetro de escala sobre a largura da função.
		}
		\label{fig:influencia-largura}
	\end{figure}
	
	\subsection{Variação da taxa de decaimento}
	
	\begin{figure}[H]
		\centering
		\includegraphics[width=0.9\textwidth]
		{influencia_decaimento.png}
		\caption{
			Influência dos parâmetros sobre a taxa de decaimento.
		}
		\label{fig:influencia-decaimento}
	\end{figure}
	
	\subsection{Observações}
	
	\textit{INSERIR ANÁLISE DOS EFEITOS DOS PARÂMETROS.}
	
	% ============================================================
	% 9. RELAÇÃO COM OS RESULTADOS EXPERIMENTAIS
	% ============================================================
	
	\section{Relação entre as Propriedades Matemáticas e os
		Resultados Experimentais}
	
	A análise matemática será posteriormente relacionada aos resultados
	experimentais obtidos para as diferentes bases.
	
	O objetivo é investigar se determinadas características matemáticas
	apresentam relação com o desempenho observado nos diferentes conjuntos
	de dados.
	
	Entre os aspectos a serem comparados estão:
	
	\begin{itemize}
		\item taxa de decaimento;
		\item concentração da influência em pontos próximos;
		\item influência de pontos distantes;
		\item suporte compacto;
		\item monotonicidade;
		\item positividade;
		\item comportamento próximo à origem;
		\item regularidade;
		\item sensibilidade aos parâmetros.
	\end{itemize}
	
	\subsection{Comparação com a acurácia}
	
	\begin{figure}[H]
		\centering
		\includegraphics[width=0.9\textwidth]
		{relacao_acuracia.png}
		\caption{
			Relação entre características matemáticas das funções e
			desempenho experimental.
		}
		\label{fig:relacao-acuracia}
	\end{figure}
	
	\subsection{Comparação com F1-score}
	
	\begin{figure}[H]
		\centering
		\includegraphics[width=0.9\textwidth]
		{relacao_f1.png}
		\caption{
			Relação entre características matemáticas das funções e
			F1-score.
		}
		\label{fig:relacao-f1}
	\end{figure}
	
	% ============================================================
	% 10. DISCUSSÃO
	% ============================================================
	
	\section{Discussão}
	
	Os resultados da análise matemática serão discutidos conjuntamente com
	os resultados experimentais.
	
	\subsection{Influência da distância}
	
	\textit{INSERIR DISCUSSÃO.}
	
	\subsection{Influência das propriedades matemáticas}
	
	\textit{INSERIR DISCUSSÃO.}
	
	\subsection{Comportamento entre diferentes conjuntos de dados}
	
	\textit{INSERIR DISCUSSÃO.}
	
	\subsection{Semelhanças e diferenças entre as bases}
	
	\textit{INSERIR DISCUSSÃO.}
	
	\subsection{Hipóteses sobre o desempenho do FEMa}
	
	\textit{INSERIR DISCUSSÃO.}
	
	% ============================================================
	% 11. CONCLUSÃO
	% ============================================================
	
	\section{Conclusão}
	
	Este documento apresentou uma proposta de análise matemática das
	21 funções de base utilizadas no algoritmo FEMa.
	
	A caracterização matemática permite identificar diferenças no
	comportamento das funções em relação à distância, seus parâmetros,
	regularidade, suporte e demais propriedades relevantes.
	
	\textit{INSERIR CONCLUSÕES OBTIDAS APÓS A ANÁLISE.}
	
	% ============================================================
	% REFERÊNCIAS
	% ============================================================
	
	\newpage
	
	\printbibliography
	
	% ============================================================
	% APÊNDICES
	% ============================================================
	
	\appendix
	
	\section{Definições Matemáticas Completas}
	
	Esta seção poderá conter as expressões matemáticas completas das
	21 funções, caso seja necessário separar as definições da análise
	principal.
	
	\section{Dados Utilizados nos Gráficos}
	
	Inserir tabelas ou informações referentes aos valores utilizados na
	geração das figuras.
	
	\section{Resultados Computacionais}
	
	Inserir resultados adicionais das verificações matemáticas realizadas
	computacionalmente.
	
\end{document}
```
